\section{Results}\label{sec:results}

\subsection{Track A: synthetic recovery}
Table~\ref{tab:track-a} reports means and standard deviations over five seeds.
When the truth is a single exponential (S1) the correctly specified single-decay
estimators (tick's exponential MLE and ADM4) are best, as they must be. Under a true
power law (S2) the sum-of-exponentials estimators---tick's and MSX---are statistically
tied (held-out deficit $0.2$ vs.\ $0.3\times10^{-3}$ nats/event; tick's solver failed on
one of five seeds). Whenever the kernel has structure outside a single family---a
delayed hump (S3), components ten milliseconds and thirty seconds apart (S4), or a sparse
ten-dimensional power-law network (S5)---MSX-auto is best on every metric: its held-out
deficit is 7--10 times smaller than the best competitor's, its branching-matrix error
is roughly halved, and it recovers the causal graph perfectly (AUC $=1$). The
exponential-only ablation isolates the contribution of Erlang phases, and the
certificate of Proposition~\ref{prop:gap} was below $10^{-4}$ nats in every MSX fit.

\begin{table}[t]\centering\scriptsize
\caption{Track A. Five seeds; dLL is the held-out log-likelihood deficit to the true model
($10^{-3}$ nats/event, lower is better); bold marks the best dLL per scenario.}
\label{tab:track-a}
\resizebox{\linewidth}{!}{\begin{tabular}{llccccc}
\toprule
Scenario & Method & $G$ err. & $\rho$ err. & kernel $L^1$ & dLL ($10^{-3}$) & AUC \\
\midrule
S1-exp & MSX-auto (ours) & 0.054 $\pm$ 0.032 & 0.020 $\pm$ 0.018 & 0.093 & 0.24 $\pm$ 0.12 & -- \\
 & MSX (ours) & 0.056 $\pm$ 0.032 & 0.019 $\pm$ 0.018 & 0.092 & 0.23 $\pm$ 0.12 & -- \\
 & MSX-R4 (ours, fixed Erlang-4) & 0.056 $\pm$ 0.035 & 0.023 $\pm$ 0.018 & 0.105 & 0.28 $\pm$ 0.15 & -- \\
 & MSX-exp (ablation: no Erlang) & 0.055 $\pm$ 0.025 & 0.026 $\pm$ 0.015 & 0.157 & 0.91 $\pm$ 0.13 & -- \\
 & tick-ExpKern(MLE, best decay) & 0.018 $\pm$ 0.007 & 0.005 $\pm$ 0.003 & 0.018 & \textbf{0.05 $\pm$ 0.05} & -- \\
 & tick-SumExpKern(MLE, AGD) & 0.065 $\pm$ 0.022 & 0.030 $\pm$ 0.013 & 0.165 & 0.92 $\pm$ 0.13 & -- \\
 & tick-ADM4 & 0.018 $\pm$ 0.007 & 0.005 $\pm$ 0.003 & 0.018 & 0.05 $\pm$ 0.05 & -- \\
 & tick-EM(nonparam) & 0.035 $\pm$ 0.015 & 0.017 $\pm$ 0.007 & 0.173 & 1.02 $\pm$ 0.07 & -- \\
 & NPHC(cumulants) & 0.247 $\pm$ 0.087 & 0.032 $\pm$ 0.024 & -- & -- & -- \\
 & tick-ConditionalLaw & 0.120 $\pm$ 0.054 & 0.017 $\pm$ 0.016 & 0.193 & 2.84 $\pm$ 1.03 & -- \\
\midrule
S2-powerlaw & MSX-auto (ours) & 0.064 $\pm$ 0.029 & 0.037 $\pm$ 0.029 & 0.108 & 0.33 $\pm$ 0.13 & -- \\
 & MSX (ours) & 0.057 $\pm$ 0.021 & 0.015 $\pm$ 0.010 & 0.114 & 0.32 $\pm$ 0.12 & -- \\
 & MSX-R4 (ours, fixed Erlang-4) & 0.062 $\pm$ 0.014 & 0.016 $\pm$ 0.006 & 0.131 & 0.45 $\pm$ 0.10 & -- \\
 & MSX-exp (ablation: no Erlang) & 0.054 $\pm$ 0.022 & 0.014 $\pm$ 0.010 & 0.097 & 0.19 $\pm$ 0.11 & -- \\
 & tick-ExpKern(MLE, best decay) & 0.299 $\pm$ 0.011 & 0.219 $\pm$ 0.007 & 0.465 & 26.05 $\pm$ 0.50 & -- \\
 & tick-SumExpKern(MLE, AGD) & 0.057 $\pm$ 0.012 & 0.016 $\pm$ 0.007 & 0.094 & \textbf{0.18 $\pm$ 0.12} & -- \\
 & tick-ADM4 & 0.299 $\pm$ 0.011 & 0.219 $\pm$ 0.007 & 0.465 & 26.05 $\pm$ 0.50 & -- \\
 & tick-EM(nonparam) & 0.057 $\pm$ 0.014 & 0.037 $\pm$ 0.008 & 0.188 & 1.73 $\pm$ 0.21 & -- \\
 & NPHC(cumulants) & 0.153 $\pm$ 0.154 & 0.041 $\pm$ 0.047 & -- & -- & -- \\
 & tick-ConditionalLaw & 0.102 $\pm$ 0.067 & 0.029 $\pm$ 0.019 & 0.291 & 4.91 $\pm$ 0.59 & -- \\
\midrule
S3-hump & MSX-auto (ours) & 0.032 $\pm$ 0.016 & 0.016 $\pm$ 0.007 & 0.080 & \textbf{0.30 $\pm$ 0.13} & -- \\
 & MSX (ours) & 0.069 $\pm$ 0.009 & 0.037 $\pm$ 0.004 & 0.347 & 5.90 $\pm$ 0.36 & -- \\
 & MSX-R4 (ours, fixed Erlang-4) & 0.038 $\pm$ 0.010 & 0.017 $\pm$ 0.005 & 0.074 & 0.31 $\pm$ 0.15 & -- \\
 & MSX-exp (ablation: no Erlang) & 0.034 $\pm$ 0.009 & 0.015 $\pm$ 0.004 & 0.416 & 11.51 $\pm$ 0.46 & -- \\
 & tick-ExpKern(MLE, best decay) & 0.166 $\pm$ 0.008 & 0.053 $\pm$ 0.004 & 0.801 & 101.47 $\pm$ 1.81 & -- \\
 & tick-SumExpKern(MLE, AGD) & 0.038 $\pm$ 0.011 & 0.018 $\pm$ 0.005 & 0.419 & 11.51 $\pm$ 0.46 & -- \\
 & tick-ADM4 & 0.166 $\pm$ 0.008 & 0.053 $\pm$ 0.004 & 0.801 & 101.47 $\pm$ 1.81 & -- \\
 & tick-EM(nonparam) & 0.024 $\pm$ 0.011 & 0.010 $\pm$ 0.004 & 0.187 & 2.24 $\pm$ 0.12 & -- \\
 & NPHC(cumulants) & 0.240 $\pm$ 0.096 & 0.045 $\pm$ 0.048 & -- & -- & -- \\
 & tick-ConditionalLaw & 0.149 $\pm$ 0.062 & 0.017 $\pm$ 0.010 & 0.319 & 10.66 $\pm$ 2.13 & -- \\
\midrule
S4-multiscale & MSX-auto (ours) & 0.084 $\pm$ 0.014 & 0.016 $\pm$ 0.006 & 0.128 & \textbf{0.32 $\pm$ 0.08} & 1.000 \\
 & MSX (ours) & 0.093 $\pm$ 0.015 & 0.021 $\pm$ 0.005 & 0.173 & 0.54 $\pm$ 0.12 & 1.000 \\
 & MSX-R4 (ours, fixed Erlang-4) & 0.101 $\pm$ 0.019 & 0.027 $\pm$ 0.007 & 0.207 & 0.70 $\pm$ 0.11 & 1.000 \\
 & MSX-exp (ablation: no Erlang) & 0.084 $\pm$ 0.014 & 0.016 $\pm$ 0.006 & 0.128 & \textbf{0.32 $\pm$ 0.08} & 1.000 \\
 & tick-ExpKern(MLE, best decay) & 0.173 $\pm$ 0.041 & 0.095 $\pm$ 0.025 & 1.635 & 711.28 $\pm$ 53.09 & 1.000 \\
 & tick-ADM4 & 0.399 $\pm$ 0.003 & 0.255 $\pm$ 0.001 & 0.410 & 3.13 $\pm$ 0.25 & 1.000 \\
 & tick-EM(nonparam) & 0.153 $\pm$ 0.012 & 0.045 $\pm$ 0.010 & 0.358 & 4.55 $\pm$ 0.29 & 1.000 \\
 & NPHC(cumulants) & 1.030 $\pm$ 0.304 & 0.211 $\pm$ 0.186 & -- & -- & 0.588 \\
 & tick-ConditionalLaw & 0.464 $\pm$ 0.286 & 0.139 $\pm$ 0.050 & 1.059 & 114.94 $\pm$ 1.20 & 0.831 \\
\midrule
S5-network10 & MSX-auto (ours) & 0.045 $\pm$ 0.004 & 0.015 $\pm$ 0.005 & 0.079 & \textbf{0.74 $\pm$ 0.10} & 1.000 \\
 & MSX (ours) & 0.111 $\pm$ 0.013 & 0.040 $\pm$ 0.013 & 0.180 & 1.19 $\pm$ 0.07 & 1.000 \\
 & MSX-R4 (ours, fixed Erlang-4) & 0.122 $\pm$ 0.008 & 0.051 $\pm$ 0.013 & 0.218 & 1.63 $\pm$ 0.09 & 1.000 \\
 & MSX-exp (ablation: no Erlang) & 0.099 $\pm$ 0.014 & 0.030 $\pm$ 0.014 & 0.145 & 0.82 $\pm$ 0.08 & 1.000 \\
 & tick-ExpKern(MLE, best decay) & 0.337 $\pm$ 0.153 & 0.103 $\pm$ 0.032 & 1.302 & 426.80 $\pm$ 368.57 & 0.977 \\
 & tick-ADM4 & 0.167 $\pm$ 0.003 & 0.096 $\pm$ 0.003 & 0.319 & 26.51 $\pm$ 0.29 & 1.000 \\
 & tick-EM(nonparam) & 0.085 $\pm$ 0.008 & 0.045 $\pm$ 0.008 & 0.280 & 5.23 $\pm$ 0.25 & 1.000 \\
 & NPHC(cumulants) & 1.348 $\pm$ 0.477 & 0.520 $\pm$ 0.372 & -- & -- & 0.672 \\
 & tick-ConditionalLaw & 0.461 $\pm$ 0.047 & 0.037 $\pm$ 0.026 & 0.641 & 78.64 $\pm$ 4.79 & 0.950 \\
\bottomrule
\end{tabular}}
\end{table}
